\section{Deletion maps and supported extensions}\label{deletion-maps-and-supported-extensions}

A deletion for \(S\subseteq I\) is an endomorphism \(e_S:G\to G\)
satisfying \(e_S(r_i)=r_i\) for \(i\in S\) and \(e_S(r_i)=1\) otherwise.
Requiring these values matters: a general retraction onto \(G_S\) need
not be a deletion.

\textbf{Proposition UA.D (Boolean deletions).} If a marked group has a
deletion \(e_S\) for every \(S\subseteq I\), then \[
e_Se_T=e_{S\cap T},\qquad \operatorname{im}e_S=\operatorname{Fix}(e_S)=G_S.
\tag{D1}
\] The marking has IP, and \[
\ker e_S=\langle\!\langle r_i:i\notin S\rangle\!\rangle_G,
\qquad G=(\ker e_S)\rtimes G_S.
\tag{D2}
\]

\emph{Proof.} The two maps in the first equality agree on every
generator. The image is generated by precisely the retained markers, and
each of those markers is fixed; this proves the image and fixed-point
statements. An element of \(G_S\cap G_T\) is fixed by
\(e_Se_T=e_{S\cap T}\), so belongs to \(G_{S\cap T}\). The reverse
inclusion is immediate.

Let \(M\) be the normal closure in (D2). Clearly \(M\leq\ker e_S\). In
\(G/M\), every omitted marker is trivial, so the image of any word
equals the image of its deletion. A word in \(\ker e_S\) is therefore
trivial in \(G/M\), proving equality. Finally, for \(g\in G\), the
factorization \(g=(g e_S(g)^{-1})e_S(g)\) has factors in \(\ker e_S\)
and \(G_S\). Their intersection is trivial because \(e_S\) fixes
\(G_S\). This is the asserted semidirect product. \(\square\)

We now assemble extensions whose extra kernel is invisible on each
proper part of its support. The groups and kernels may be infinite and
nonabelian.

\textbf{Theorem UA.SK (supported-kernel realization).} Let \(I\) be
finite and let \(\mathcal F\) be a finite family of distinct subsets of
\(I\). Let \(B=\langle b_i:i\in I\rangle\) have IP. For every
\(S\in\mathcal F\), suppose the following data are given:

\begin{enumerate}
\def\labelenumi{\arabic{enumi}.}
\tightlist
\item
  A deletion homomorphism \(p_S:B\to B_S\), fixing \(b_i\) for
  \(i\in S\) and killing \(b_i\) otherwise.
\item
  A group \(E_S=\langle e_{S,i}:i\in S\rangle\) and a marked epimorphism
  \(q_S:E_S\to B_S\).
\item
  Injectivity of \(q_S\) on \((E_S)_T=\langle e_{S,i}:i\in T\rangle\)
  for every \(T\subsetneq S\).
\end{enumerate}

Put \(N_S=\ker q_S\), and form the group \[
H=\left\{(b,(x_S)_{S\in\mathcal F})\in
B\times\prod_{S\in\mathcal F}E_S:
q_S(x_S)=p_S(b)\text{ for every }S\right\}.
\tag{K1}
\] For \(i\in I\), define \(h_i\) to have base coordinate \(b_i\) and
\(S\)-coordinate \(e_{S,i}\) when \(i\in S\), and \(1\) otherwise. These
elements generate all of \(H\). If \(\pi:H\to B\) is the base projection
and \(H_L=\langle h_i:i\in L\rangle\), then, for every \(L\subseteq I\),
\[
\pi(H_L)=B_L,\qquad
\ker\pi=\prod_{S\in\mathcal F}N_S,\qquad
\ker\pi\cap H_L=\prod_{\substack{S\in\mathcal F\\S\subseteq L}}N_S.
\tag{K2}
\] The factors on the right are embedded as pure coordinate subgroups.
The marking of \(H\) has IP. If \(B\) and all \(E_S\) are finite, then
\[
|H_L|=|B_L|\prod_{\substack{S\in\mathcal F\\S\subseteq L}}|N_S|.
\tag{K3}
\]

\emph{Proof.} The compatibility equations in (K1) are preserved by
products and inverses, so \(H\) is a subgroup of the displayed direct
product. The prescribed tuple \(h_i\) satisfies those equations by the
literal deletion rule. Let \(H'\) be the subgroup generated by these
tuples. Its base image is \(B\), and the projection of
\(H'_S=\langle h_i:i\in S\rangle\) onto \(E_S\) is onto.

We first prove, by induction on inclusion among the finitely many
supports, that every pure coordinate element of \(N_S\) belongs to
\(H'_S\). Fix \(n\in N_S\). Choose a word on the markers indexed by
\(S\) which evaluates to \(n\) in \(E_S\), and evaluate the same word in
\(H'\), obtaining \(w\in H'_S\). Its base coordinate lies in \(B_S\) and
is \(q_S(n)=1\).

Consider its coordinate at another support \(T\). Because the word uses
only markers from \(S\), that coordinate lies in \((E_T)_{S\cap T}\).
Its image under \(q_T\) is \(p_T(1)=1\). If \(T\not\subseteq S\), then
\(S\cap T\subsetneq T\), and the proper-parabolic injectivity hypothesis
forces this coordinate to be \(1\). The only possible nonidentity
coordinates of \(w\) are consequently at supports \(T\subseteq S\). Its
\(S\)-coordinate is the prescribed \(n\), while every coordinate at a
proper support \(T\subsetneq S\) belongs to \(N_T\).

For each such proper support, the induction hypothesis supplies the pure
coordinate element having exactly that value in \(N_T\), as an element
of \(H'_T\leq H'_S\). Multiply \(w\) by the inverses of those elements.
Distinct pure coordinate subgroups commute in the ambient direct
product, so these operations remove the unwanted coordinates without
changing the \(S\)-coordinate. The result is the pure element \(n\) in
coordinate \(S\), and remains in \(H'_S\). A minimal support has no
smaller coordinates to remove, so this argument includes the initial
step. No commutativity within an individual kernel has been used.

The kernel of \(\pi\) in the full group \(H\) is exactly the direct
product of the \(N_S\), by (K1). The induction places that entire kernel
in \(H'\). Since \(H'\) already maps onto \(B\), every \(h\in H\) has
the same base coordinate as some \(h'\in H'\); then
\(h(h')^{-1}\in\ker\pi\leq H'\). Thus \(H'=H\).

The image equality \(\pi(H_L)=B_L\) follows on generators. If
\(x\in H_L\cap\ker\pi\), its \(S\)-coordinate lies in
\((E_S)_{L\cap S}\) and in \(N_S\). This coordinate is trivial whenever
\(S\not\subseteq L\), again by proper-parabolic injectivity. Conversely,
when \(S\subseteq L\), the isolation argument placed the entire pure
\(N_S\) in \(H_S\leq H_L\). This proves all three assertions in (K2).

The coordinate products in (K2) satisfy \[
\left(\prod_{S\subseteq L}N_S\right)\cap
\left(\prod_{S\subseteq M}N_S\right)
=\prod_{S\subseteq L\cap M}N_S.
\] Theorem UA.K now proves IP in \(H\). More explicitly, an element of
\(H_L\cap H_M\) has base in \(B_{L\cap M}\); subtracting a lift from
\(H_{L\cap M}\) leaves an element in both displayed kernel products,
hence in the kernel product for \(L\cap M\). This also shows exactly
where IP of the base is used. Finally (K3) is the order formula for the
surjection \(H_L\to B_L\) with the kernel just computed. \(\square\)

\textbf{Corollary UA.SK.Local (recognition on an antichain).} Under
UA.SK, if \(\mathcal F\) is an antichain under inclusion, then the
coordinate projection \(H_S\to E_S\) is a marked isomorphism for every
\(S\in\mathcal F\).

\emph{Proof.} It is onto on generators. If an element maps to \(1\), its
base coordinate also vanishes, since the base value of a word on \(S\)
is the \(q_S\)-image of its \(S\)-coordinate. Formula (K2) says that the
remaining kernel in \(H_S\) has only the factor \(N_S\): the antichain
excludes every other support contained in \(S\). The \(S\)-coordinate
detects that factor, so the element is \(1\). \(\square\)

\textbf{Corollary UA.SK.Delete (derived deletions).} If, in addition,
\(B\) has deletions for every subset of \(I\), then every \(E_S\) and
\(H\) has all its Boolean deletions.

\emph{Proof.} For \(T\subsetneq S\), the restriction
\(q_S:(E_S)_T\to B_T\) is a marked isomorphism: it is onto on generators
and injective by hypothesis. Compose \(q_S\), the deletion
\(B_S\to B_T\), and the inverse of this isomorphism. This gives a
homomorphism \(E_S\to(E_S)_T\) that fixes exactly the \(T\)-markers and
kills the others. For \(T=S\) use the identity. Proposition UA.D
supplies the Boolean identities.

For \(L\subseteq I\), apply the base deletion \(e_L\) and, in coordinate
\(S\), the local deletion onto \(S\cap L\). The two sides of each
compatibility equation agree on the local markers; the base deletions
commute by UA.D. Thus the coordinatewise operation preserves (K1). On
\(h_i\) it fixes \(h_i\) if \(i\in L\), and kills it otherwise. It is
the required deletion of \(H\). \(\square\)

\textbf{Corollary UA.SK.Central (lower-central kernel).} Under UA.SK,
suppose \(c\geq1\), \(\gamma_c(B)=1\), and \(N_S=\gamma_c(E_S)\) for
every support. Then \[
\gamma_c(H)=\ker\pi.
\tag{K4}
\]

\emph{Proof.} Applying \(\pi\) to a commutator of weight \(c\) gives
\(1\), so \(\gamma_c(H)\leq\ker\pi\). For the converse, repeat the
inclusion induction in the proof of UA.SK. The projection \(H_S\to E_S\)
is onto; therefore its image on \(\gamma_c(H_S)\) is
\(\gamma_c(E_S)=N_S\). Given \(n\in N_S\), choose a lift
\(w\in\gamma_c(H_S)\). Its base is \(1\), since \(\gamma_c(B)=1\). The
same proper-support argument makes all coordinates at
\(T\not\subseteq S\) trivial. The remaining proper-support coordinates
lie in \(N_T\); by induction their pure copies belong to
\(\gamma_c(H)\). Removing them leaves the pure \(n\), still in
\(\gamma_c(H)\). All pure kernels are obtained, proving (K4).
\(\square\)

\textbf{Corollary UA.SK.Separator (factorization across a separator).}
Under UA.SK, let \(I=L\cup M\). Suppose \(B=B_LB_M\) as a product of
sets, and every support in \(\mathcal F\) is contained in \(L\) or in
\(M\). Then \(H=H_LH_M\).

\emph{Proof.} For \(h\in H\), factor \(\pi(h)=b_Lb_M\) and choose lifts
\(u_L\in H_L\), \(u_M\in H_M\). Then \(h=n u_Lu_M\) for some
\(n\in\ker\pi\). Assign every pure coordinate factor of \(n\) to one of
the two sides containing its support. Formula (K2) gives \(n=n_Ln_M\),
with \(n_L\in H_L\), \(n_M\in H_M\). Each pure \(N_S\) is normal in
\(H\), because conjugation in the ambient product acts coordinatewise
and \(N_S\unlhd E_S\). Hence the product of the assigned \(M\)-factors
is normal in \(H\). Moving \(n_M\) past \(u_L\) gives \[
h=(n_Lu_L)(u_L^{-1}n_Mu_L\,u_M)\in H_LH_M.
\] The opposite inclusion is automatic. \(\square\)

The hypotheses of UA.SK do not identify \(H\) with a universal amalgam
of the \(E_S\). Additional cross-relations and a
presentation-recognition argument may be necessary. Nor can injectivity
on every proper parabolic be replaced by injectivity on individual
markers alone.

\textbf{Proposition UA.Fiber (two compatible split factors).} Let \(A\)
and \(B\) be marked groups on finite label sets \(L\) and \(M\), and put
\(J=L\cap M\). Suppose their \(J\)-parabolics are identified with the
same marked group \(C\), and that \(p_A:A\to C\) and \(p_B:B\to C\) are
deletions. Then \[
A\times_C B\cong(\ker p_A\times\ker p_B)\rtimes C.
\tag{F1}
\] This full fiber product is generated by the tuples \((a_i,1)\) for
\(i\in L\setminus J\), \((1,b_i)\) for \(i\in M\setminus J\), and
\((c_i,c_i)\) for \(i\in J\). The two original groups embed through
\(a\mapsto(a,p_A(a))\) and \(b\mapsto(p_B(b),b)\). If finite, its order
is \(|A||B|/|C|\).

\emph{Proof.} Every compatible pair has a unique form \((n_Ac,n_Bc)\),
giving the split decomposition and the order formula. The stated copies
of \(A\) and \(B\) are generated by subsets of the stated tuple marking.
They contain \((n_A,1)\) and \((1,n_B)\), respectively, while either
contains the diagonal copy of \(C\). These subgroups generate the split
product, proving the assertion. The embedding formulas are injective
because their first and second projections, respectively, are
identities. \(\square\)
